% Part: first-order-logic
% Chapter: completeness
% Section: outline

\documentclass[../../../include/open-logic-section]{subfiles}

\begin{document}

\iftag{FOL}
      {\olfileid[pt]{fol}{com}{out}}
      {\olfileid[pt]{pl}{com}{out}}

\olsection{Esboço da Prova}

A prova do teorema da completude é um pouco complexa e, em uma primeira
leitura, é fácil se perder.  Então, vamos esboçar a prova.  O
primeiro passo é uma mudança de perspectiva, que nos permite enxergar
um caminho para uma prova.  Quando a completude é pensada como
``sempre que $\Gamma \Entails !A$ então $\Gamma \Proves !A$'', pode
ser difícil até mesmo ter uma ideia: pois, para mostrar que $\Gamma
\Proves !A$, temos que encontrar !!a{derivation}, e não parece que a
hipótese de que $\Gamma \Entails !A$ nos ajude nisso de alguma forma.
Para alguns sistemas de prova, é possível construir diretamente
!!a{derivation}, mas adotaremos uma abordagem um pouco diferente.  A
mudança de perspectiva necessária é esta: a completude também pode ser
formulada como: ``se $\Gamma$ é consistente, então é satisfatível''.
Talvez possamos usar a informação em~$\Gamma$ juntamente com a hipótese
de que ela é consistente para construir
\iftag{FOL}{!!a{structure}}{!!a{valuation}} que satisfaça toda
\iftag{FOL}{!!{sentence}}{!!{formula}} em~$\Gamma$.  Afinal, sabemos
que tipo de \iftag{FOL}{!!{structure}}{!!{valuation}} estamos
procurando: uma que seja como $\Gamma$ a descreve!{}

Se $\Gamma$ contém apenas \iftag{FOL}{!!{sentence}s
  atômicas}{!!{propositional variable}s}, é fácil construir
um modelo para ele.\iftag{FOL}{ Suponha que as !!{sentence}s atômicas
  sejam todas da forma $\Atom{P}{a_1,\dots,a_n}$ onde os $a_i$ são
  !!{constant}s.}{} Tudo o que temos a fazer é apresentar
\iftag{FOL}{%
  !!a{domain}~$\Domain{M}$ e uma atribuição para~$P$ de modo que
  $\Sat{M}{\Atom{P}{a_1,\ldots,a_n}}$. Mas isso não é muito difícil:
  faça $\Domain{M} = \Nat$, $\Assign{\Obj c_i}{M} = i$ e, para cada
  $\Atom{P}{a_1,\ldots,a_n} \in \Gamma$, coloque a tupla $\tuple{k_1,
    \dots, k_n}$ em $\Assign{P}{M}$, onde $k_i$ é o índice do
  símbolo de constante~$a_i$ (isto é, $a_i \ident \Obj c_{k_i}$).}{%
  !!a{valuation}~$\pAssign{v}$ tal que $\pSat{v}{p}$ para todo $p \in
  \Gamma$. Bem, faça $\pAssign{v}(p) = \True$ se e somente se $p \in \Gamma$.}

Agora suponha que $\Gamma$ contenha alguma !!{formula}~$\lnot !B$, com $!B$
atômica.  Poderíamos temer que a construção de
$\iftag{FOL}{\Struct{M}}{\pAssign{v}}$ interfira na possibilidade
de tornar $\lnot !B$ verdadeira.  Mas é aqui que entra a consistência
de~$\Gamma$: se $\lnot !B \in \Gamma$, então $!B \notin
\Gamma$, caso contrário $\Gamma$ seria inconsistente.  E se $!B \notin
\Gamma$, então, de acordo com nossa construção
de~$\iftag{FOL}{\Struct{M}}{\pAssign{v}}$,
$\iftag{FOL}{\Sat/{M}}{\pSat/{v}}{!B}$, logo
$\iftag{FOL}{\Sat{M}}{\pSat{v}}{\lnot !B}$.  Até aqui, tudo bem.

E se $\Gamma$ contiver fórmulas complexas, não atômicas? Digamos que ele
contenha $!A \land !B$. Para tornar isso verdadeiro, devemos proceder como se
tanto $!A$ quanto $!B$ estivessem em~$\Gamma$.  E se $!A \lor !B \in \Gamma$,
então teremos que tornar pelo menos um deles verdadeiro, isto é, proceder como
se um deles estivesse em~$\Gamma$.

Isso sugere a seguinte ideia: adicionamos !!{formula}s adicionais
a~$\Gamma$ de modo a (a)~manter o conjunto resultante consistente e
(b)~garantir que, para toda !!{sentence}~$!A$ atômica possível,
ou $!A$ está no conjunto resultante, ou $\lnot !A$ está, e (c)~tais que,
sempre que $!A \land !B$ está no conjunto, então tanto $!A$ quanto $!B$
também estão; se $!A \lor !B$ está no conjunto, pelo menos um dentre $!A$
ou $!B$ também está, etc.  Continuamos fazendo isso (potencialmente para
sempre).  Chame o conjunto de todas as
!!{formula}s assim adicionadas de~$\Gamma^*$.  Então a nossa construção acima
nos forneceria \iftag{FOL}{!!a{structure}~$\Struct{M}$}{!!a{valuation}~$\pAssign{v}$}
para a qual poderíamos provar, por indução, que ela satisfaz todas as sentenças
em~$\Gamma^*$ e, portanto, também todas as sentenças em~$\Gamma$,
já que $\Gamma \subseteq \Gamma^*$.  Acontece que garantir (a)
e~(b) é suficiente. Um conjunto de sentenças para o qual (b) vale é chamado de
\emph{completo}. Assim, nossa tarefa será estender o conjunto consistente
~$\Gamma$ a um conjunto consistente e completo~$\Gamma^*$.

\iftag{FOL}{%
Há um problema neste plano: se $\lexists[x][!A(x)] \in \Gamma$,
esperaríamos poder escolher algum !!{constant}~$c$ e adicionar $!A(c)$
neste processo.  Mas como sabemos que sempre podemos fazer isso?  Talvez
tenhamos apenas alguns !!{constant}s em nossa linguagem, e, para cada um
deles, tenhamos $\lnot !A(c) \in \Gamma$.  Também não podemos adicionar $!A(c)$,
já que isso tornaria o conjunto inconsistente, e não saberíamos
se $\Struct{M}$ tem que tornar $!A(c)$ ou $\lnot !A(c)$ verdadeira.
Além disso, pode acontecer de $\Gamma$ conter apenas sentenças em uma
linguagem que não tem símbolo de constante algum (por exemplo, a linguagem da
teoria dos conjuntos).

A solução para este problema é simplesmente adicionar infinitas
constantes no começo, além de sentenças que as conectam com os
quantificadores da maneira correta.  (Claro, temos que verificar que isso
não pode introduzir uma inconsistência.)

Nossa construção original funciona bem se tivermos apenas !!{constant}s nas
sentenças atômicas.  Mas a linguagem também pode conter
!!{function}s.  Nesse caso, pode ser complicado encontrar as funções certas
em~$\Nat$ para atribuir a esses !!{function}s de modo a fazer
tudo funcionar.  Então, aqui vai outro truque: em vez de usar $i$ para
interpretar $\Obj c_i$, basta tomar o próprio conjunto de !!{constant}s como o
domínio.  Então $\Struct M$ pode atribuir cada !!{constant} a si mesmo:
$\Assign{\Obj c_i}{M} = \Obj c_i$.  Mas por que não ir até o fim: seja
$\Domain{M}$ o conjunto de todos os \emph{termos} da linguagem!{} Se fizermos isso,
há uma atribuição óbvia de funções (que tomam termos como
argumentos e têm termos como valores) aos !!{function}s: atribuímos ao
!!{function}~$\Obj f^n_i$ a função que, dados $n$ termos $t_1$,
\dots,~$t_n$ como entrada, produz o termo $\Obj f^n_i(t_1, \dots, t_n)$
como valor.

A última peça do quebra-cabeça é o que fazer com~$\eq$.  O
!!{predicate}~$\eq$ tem uma interpretação fixa: $\Sat{M}{\eq[t][t']}$
se e somente se $\Value{t}{M} = \Value{t'}{M}$.  Agora, se ajustarmos as coisas de modo que o
!!{value} de um termo~$t$ seja o próprio $t$, então esta !!{structure}
não tornará verdadeira \emph{nenhuma} sentença da forma $\eq[t][t']$
a menos que $t$ e $t'$ sejam um único e mesmo termo.  E, claro,
isso é um problema, já que basicamente toda teoria interessante em
uma linguagem com !!{function}s terá como teoremas sentenças
$\eq[t][t']$ onde $t$ e $t'$ não são o mesmo termo (por exemplo, em
teorias da aritmética: $\eq[(\Obj 0+ \Obj 0)][\Obj 0]$).  Para resolver
este problema, mudamos o domínio de~$\Struct M$: em vez de usar termos como
os objetos em~$\Domain{M}$, usamos conjuntos de termos, e cada conjunto é tal
que contém todos aqueles termos que as sentenças em~$\Gamma$
exigem que sejam iguais.  Assim, por exemplo, se $\Gamma$ é uma teoria da aritmética,
um desses conjuntos conterá: $\Obj 0$, $(\Obj 0 + \Obj 0)$, $(\Obj
0 \times \Obj 0)$, etc.  Este será o conjunto que atribuímos a $\Obj 0$,
e acontecerá que este conjunto também é o valor de todos os termos
nele, por exemplo, também de $(\Obj 0 + \Obj 0)$.  Portanto, a sentença
$\eq[(\Obj 0+ \Obj 0)][\Obj 0]$ será verdadeira nesta
!!{structure} revisada.}{}

Então, aqui está o que faremos. Primeiro investigamos as propriedades dos
conjuntos consistentes !!{complete}s; em particular, provamos que
um conjunto consistente !!{complete} contém $!A \land !B$ se e somente se contém
tanto $!A$ quanto~$!B$, contém $!A \lor !B$ se e somente se contém pelo menos um deles,
etc. (\olref[ccs]{prop:ccs}).\iftag{FOL}{ Em seguida, definimos e
  investigamos conjuntos ``saturados'' de sentenças. Um conjunto saturado é aquele
  que contém condicionais que ligam cada !!{sentence} quantificada
  às suas instâncias (\olref[hen]{defn:henkin-exp}). Mostramos que qualquer
  conjunto consistente~$\Gamma$ sempre pode ser estendido a um conjunto saturado
  ~$\Gamma'$ (\olref[hen]{lem:henkin}). Se um conjunto é consistente,
  saturado e !!{complete}, ele também tem a propriedade de
  conter \iftag{prvEx}{$\lexists[x][!A(x)]$ se e somente se contém $!A(t)$
    para algum termo fechado~$t$}{}\iftag{defEx,defAll}{}{ e
  }\iftag{prvAll}{$\lforall[x][!A(x)]$ se e somente se contém $!A(t)$ para todo
    termo fechado~$t$}{} (\olref[hen]{prop:saturated-instances}).}{}
Em seguida, tomaremos o conjunto consistente\iftag{FOL}{ saturado}{}
~$\iftag{FOL}{\Gamma'}{\Gamma}$ e mostraremos que ele pode ser estendido
a um conjunto \iftag{FOL}{saturado, consistente e}{consistente e}
!!{complete}~$\Gamma^*$ (\olref[lin]{lem:lindenbaum}). Este conjunto
$\Gamma^*$ é o que usaremos para definir \iftag{FOL}{nosso modelo de termos~$\Struct M(\Gamma^*)$}{nossa !!{valuation}~$\pAssign v(\Gamma^*)$}.
\iftag{FOL}{O modelo de termos tem o conjunto de termos fechados como seu domínio,
  e a interpretação de seus !!{predicate}s é dada pelas !!{sentence}s
  atômicas}{A valoração é determinada pelas !!{propositional
    variable}s} em~$\Gamma^*$ (\olref[mod]{defn:termmodel}).  Usaremos
as propriedades dos conjuntos consistentes completos\iftag{FOL}{ e saturados}{}
para mostrar que, de fato,
$\iftag{FOL}{\Sat{M(\Gamma^*)}}{\pSat{v(\Gamma^*)}}{!A}$ se e somente se $!A \in
\Gamma^*$ (\olref[mod]{lem:truth}) e, portanto, em particular,
$\iftag{FOL}{\Sat{M(\Gamma^*)}}{\pSat{v(\Gamma^*)}}{\Gamma}$.
\iftag{FOL}{Por fim, consideraremos como definir um modelo de termos se
  $\Gamma$ contém~$\eq$ também
  (\olref[ide]{defn:term-model-factor}) e mostraremos que ele
  satisfaz~$\Gamma^*$ (\olref[ide]{lem:truth}).}{}

\end{document}
